\section{Result and comparison contract}
\label{sec:assessment}

The proposed presentation has six primitive inference patterns and the
numerical term constructors $0$ and $R$, in addition to variables. Successor
is a derived notation. There is one explicit computation schema. Its
ordinary realization is bi-interpretable with full first-order intuitionistic
primitive recursive arithmetic, with the entire numerical domain in both
directions and with identity graph $x=y$ for both numerical composites.

The reduction is precise but has a cost: the zero computation becomes a
canonical-syntax convention, and substitution must enforce that convention.
We count this normalization in the complete implementation. A reduction from
seven named proof rules to six is therefore not, on its own, a proof of a
smaller complete checker.

\begin{center}
\small
\begin{tabular}{lll}
\toprule
Feature & Current scalar presentation & Canonical $R$ presentation\\
\midrule
Numerical constructors & $0,S,I$ & $0,R$\\
Explicit computation schemas & Two & One\\
Primitive inference patterns & Seven & Six\\
Term equality in the checker & Literal syntax & Canonical syntax\\
Substitution & Capture avoiding & Capture avoiding, then normalize\\
Numerical interpretation domain & Whole & Whole\\
Numerical composite graph & $x=y$ & $x=y$\\
\bottomrule
\end{tabular}
\end{center}

The comparison is pinned to ProveIt commit
\texttt{8eb71c1e570e07ef45ead686628e77f18ef34798}, inspected on
12 November 2026, UTC. The current seven-rule scalar calculus is
\texttt{MPrf} in \texttt{NativeMinimal.lean}. The separate crossed calculus
\texttt{RPrf} is in \texttt{NativeReversal.lean}. Their equivalences with the
native scalar core are already formalized upstream. The present canonical
$R$ reduction has a written proof and executable implementation tests; it is
not a new Lean or Rocq formalization.

The repository's smallest compact profiles are a different comparison
category. The selected byte profile has 522 normalized source bytes, but its
bi-interpretation theorem concerns a guarded public language, with broader
auxiliary proof syntax. Its compact program does not itself validate that
public grammar. Our strict canonical-$R$ checker validates its full public
language. We therefore report two separate comparisons: matching strict
implementations of the mathematical calculi, and small implementation
changes to the existing compact calculus. No compact size is attributed to
the new six-rule mathematics.

The successor-seeded recursor, its predecessor construction, the scalar
arithmetic bridge, crossed transport, and witness-directed introduction are
inherited. The mathematical change developed here is canonicalization of the
successor-seeded syntax and the proof that it preserves substitution,
derivability, and both interpretation composites. A separate strengthening
of the scalar self-seeded obstruction is proved in
Section~\ref{sec:scalar-obstruction}.

